\section{Conclusion}
\label{sec:conclusion}

This project analysed two data sets, a lab with HTTP and SYN floods and a public capture of an infected Windows workstation, and took both through the complete security-monitoring process (Table~\ref{tab:nsm-process}).

\begin{table}[htbp]
  \centering
  \caption{The security-monitoring process applied in this project.}
  \label{tab:nsm-process}
  \footnotesize
  \begin{tabularx}{\linewidth}{@{}lY@{}}
    \toprule
    \textbf{Stage} & \textbf{What the project did} \\
    \midrule
    Network data & Captured lab traffic; a 53-minute malware capture, hashed before use \\
    Collection & Zeek and Suricata logs loaded into Elasticsearch and Kibana; probe, resource and nginx data for the lab \\
    Analysis & Packet timing, flow and request-level analysis, security dashboard \\
    Security indicator & Request floods and half-open handshakes; bare-IP polling every 5.3 s, self-signed TLS on port 2917 \\
    Investigation & Host DESKTOP-L8C5GSJ and user shutchenson identified; infection chain reconstructed \\
    Detection & Indicators turned into rules; tuning cut alerts from 619 to 18 \\
    Decision & Compromise confirmed with high confidence; floods caused degradation, not compromise \\
    Response & Isolate and reimage the host, block C2 infrastructure, feed new rules back \\
    \bottomrule
  \end{tabularx}
\end{table}

Three lessons stand out. Suspicious behaviour was recognised by combinations of features judged against the host's role, not by a single threshold. Every finding needed a second, independent source before it was accepted. And indicators were short-lived, because the C2 servers were set up hours before use, so behavioural rules such as request-level polling detection gave the durable results.

Network security monitoring is therefore not the collection of packets and logs. Its value comes from turning that data into reliable security evidence that supports detection, investigation and incident response.
